% RESEARCH_DRAFT_SHA256: 774172c507fb54fc557d96eb2e25f7efe2a3351d0ae33c167d0ea23eb12bb3e3
% Layout and source-reference wrappers expanded for the web renderer.
\documentclass[11pt]{article}
\usepackage[margin=25mm]{geometry}
\usepackage{amsthm}
\usepackage{fontspec,unicode-math}
\setmainfont{Latin Modern Roman}
\setmathfont{Latin Modern Math}
\usepackage{amsmath,amssymb,mathtools,xurl,hyperref}
\hypersetup{colorlinks=true,urlcolor=blue,linkcolor=blue}
\setcounter{secnumdepth}{0}
\setlength{\parindent}{0pt}
\setlength{\parskip}{5pt}
\setlength{\emergencystretch}{3em}
\newcommand{\sref}[2]{\hyperlink{source:#1:#2}{[S#2]}}
\newcommand{\cataloguescope}[1]{\paragraph{Source record.}#1}
\newcommand{\floor}[1]{\left\lfloor#1\right\rfloor}
\newtheorem{proposition}{Proposition}
\newtheorem{lemma}[proposition]{Lemma}
\begin{document}
\paragraph{Research manuscript.} OpenAI \texttt{gpt-6-astra}, reasoning \texttt{ultra}; 11 October 2026.
\paragraph{Preserved source body.} The complete catalogue mathematical text and sources follow with layout and source links expanded, including source qualifications and its historical preparation notes. New research begins at the section headed Research attempt.

% UnsolvedMath ID 30005827; source catalogue code OWR-14298166-003
\section{Average-Case Complexity of Clique on Sparse Graphs}\label{30005827}
{\small UnsolvedMath ID 30005827 · Computer Science}
\subsection{Definitions and mathematical statement}

\paragraph{Shared notation.}Standard number systems are $\mathbb N,\mathbb Z,\mathbb Q,\mathbb R,\mathbb C$. Specialized notation and index ranges follow the cited source and the definitions below.


An Erdos-Renyi graph $G(n,p)$ has independent edges, each present with probability $p$. A $k$-clique is a $k$-vertex set with every possible joining edge. Average-case complexity is measured over this graph distribution, separately from worst-case clique hardness.
\paragraph{Record scope and qualifications.}
The catalogue title is retained as metadata. Named certificates, contracts, bounds or frameworks in that title are not asserted as results of the original source extract.
The exact threshold regime, success probability, parameter scaling, and machine/access model are not fully quantified in the supplied record.

Source-record limitation: The cleaned modest variant switches from k-cliques to n^epsilon-cliques while retaining a runtime parameter k, without defining their relation or the distributional success criterion. The relationship between these parameters is absent from the supplied extraction.
\paragraph{Statement recorded in the supplied source.}
For Erdős–Rényi graphs with edge probability near the threshold for containing a $k$-clique, does deciding $k$-Clique require time $n^{\Omega(k)}$? More modestly, is there an $n^{o(k)}$-time algorithm for deciding whether such a graph contains a clique of size $n^\varepsilon$? \href{https://doi.org/10.4171/owr/2024/15}{[S2]}
\paragraph{Source status.}
The supplied catalogue labels this source record open; that metadata was checked against the supplied snapshot on 2026-10-11. The label is preserved for provenance and is not an independent verification of present open status or mathematical correctness. The catalogue literature assessment is dated 2026-08-22.
\subsection{Short English statement}
Determine the average-case cost of finding or deciding cliques near their random-graph threshold.
\subsection{Sources}
\begin{itemize}
\item[S1] Ulam AI and the UnsolvedMath Contributors, UnsolvedMath: A Curated Collection of Open Mathematics Problems, record 30005827 (OWR-14298166-003). Catalogue curation and metadata: CC BY 4.0. The source extract is retained; definitions, English summaries and qualifications are editorial additions.
\url{https://huggingface.co/datasets/ulamai/UnsolvedMath}.
\item[S2] Proof Complexity and Beyond (2024), Oberwolfach Reports, DOI 10.4171/owr/2024/15, pp. 871--934
\url{https://doi.org/10.4171/owr/2024/15}.
\item[S3] de Rezende, Potechin, and Risse, Clique Is Hard on Average
\url{https://arxiv.org/abs/2404.16722}.
\item[S4] Oberwolfach Report 15/2024
\url{https://oa.tib.eu/renate/bitstreams/b0ef46dc-1bdf-45b7-a4f5-ca0a9e60a7d6/download}.
\item Source-specific Oberwolfach Reports licences, including CC BY-SA 4.0 where applicable, are retained; see THIRD\_PARTY\_NOTICES.md.
\end{itemize}
\paragraph{Preparation and review.}
The original source and UnsolvedMath contributors retain attribution for the source extract. Codex (GPT-6; exact serving revision not exposed) assisted with definitions, English summaries, scope notes and formatting. No new human mathematical authorship or independently verified solution is asserted. Permanent source ID 30005827 and code OWR-14298166-003 checked against the complete supplied registry and current repository. Distinct upstream records, including repeated questions and broad or damaged extracts, are retained for complete coverage of the source CS category. A repeated record is not a new independent conjecture. Complete mathematical-equivalence and current-literature review has not been performed. Preparation and final review follow contribution prompts C and D, with source gaps disclosed. Structural and rendering checks do not establish a new theorem or open status.
\label{end:30005827}

\section{Research attempt}
\subsection{Literature, source context and scope}
The full original catalogue body and source qualifications are preserved above. The relevant primary workshop passages were read in the cached publisher report. Historical claims and subsequent attributed results are distinguished from the mathematical work below. A restricted example, standard lemma or known theorem is not counted as new progress on the entire original target. Missing assumptions and repeated or rhetorical extracts are not silently turned into different solved problems.

\subsection{A distributional promise can make the constant answer succeed}
The primary source discusses $k$-clique near a random-graph threshold and bounded-coefficient proof lower bounds in specified parameter ranges. The short catalogue sentence does not fix how $k,p,n$ scale or whether algorithms may err on a small probability of inputs. We test a concrete bounded-error distributional interpretation, rather than silently importing a refutation requirement.
Let $k\ge3$ be fixed, $n\ge k$, and $G\sim G(n,p)$. A decision algorithm always answering ``no $k$-clique'' has no running-time dependence on the graph contents.
\begin{proposition}
Its error probability is at most $\binom nk p^{\binom k2}$. In particular, for $p=n^{-2/(k-1)}$ it is at most $1/k!\le1/6$. For $p=n^{-2/(k-1)}/\log n$, the error tends to zero as $n\to\infty$.
\end{proposition}
\begin{proof}
For each $k$-vertex subset the probability that all its $\binom k2$ edges occur is $p^{\binom k2}$, by independence. Summing these probabilities bounds the event that at least one subset is a clique. At the first stated $p$, this factor is $n^{-k}$ and $\binom nk\le n^k/k!$. At the second $p$ there is the additional factor $(\log n)^{-\binom k2}$, tending to zero.
\end{proof}
Consequently a universal lower bound for every bounded-error decision algorithm on every distribution called ``near threshold'' cannot be inferred from a proof-complexity lower bound. The constant algorithm has no certificate and can be wrong on all positive instances. A zero-error algorithm, a certifying algorithm, a balanced YES/NO experiment, growing $k$, or worst-case correctness poses a different problem. This elementary calculation does not contradict the source's specified bounded-coefficient Sherali--Adams lower bounds.
\subsection{Verification and remaining gaps}
Small random-graph spaces were exhaustively enumerated with rational $p$ to compare the exact clique probability and the first-moment bound. The threshold-scale inequality uses the exact exponent cancellation above, not those finite examples. The unspecified scaling in the original statement remains unspecified; no fine-grained lower bound is proved, and no claim is made for a correctly fixed hard average-case promise.

\section{References and source attribution}
Primary workshop report: \url{https://doi.org/10.4171/owr/2024/15}. Relevant passages read: 2024/15 pp.890--891, random-graph clique regimes and bounded-coefficient Sherali--Adams discussion.

All retrievals were checked on 11 October 2026. Original question attribution remains in the preserved Sources section. External theorems are attributed and not claimed as new discoveries or fully verified merely by citation.


\paragraph{Integration review.} Codex (GPT-6; exact serving revision not exposed) reviewed the written arguments and expanded layout and source-reference wrappers. Repository metadata and the production web renderer were checked using Python 3.13 and JavaScript with MathJax 3.2.2. This model-assisted review is not independent human proof verification; no PDF compilation or formal proof checking is asserted.

\section{Completion Estimate}
\noindent\textbf{COMPLETION: 0\%}
No substantiated new progress toward the entire original target is claimed. The worked results and failed route are delimited above; a known result, illustrative model, or auxiliary obstruction is not counted as original-target completion.

\section{Final status and attribution}
\textbf{Attempt status: unresolved.} This describes this research attempt, not a new ruling on the source catalogue's status or on external literature claims.
The mathematics and this self-audit were generated by OpenAI \texttt{gpt-6-astra}, reasoning \texttt{ultra}, on 11 October 2026. Codex (GPT-6; exact serving revision not exposed) orchestrated the batch, read and reviewed its written mathematical arguments, and contributed formatting/repository review. No human mathematical authorship, independent human proof checking, or formal verification is claimed. The historical preparation paragraphs in the reproduced catalogue body concern that body only.

Contribution prompt A was applied to the complete definition and supplied source record, followed by the adversarial checks of prompt B. Local reproduction files are \texttt{batch025/checks.py}, \texttt{computations.json}, and the per-ID audit JSON. They were not public links at preparation time. Numerical checks are finite and do not prove unrestricted results. No LaTeX engine was found locally; compilation was not performed by this research agent.

\end{document}
